\section{Finite data, coding, and local contracts}
\label{sec:finite-data}

\subsection{One numerical sort is enough}

All finite data used by the framework are natural numbers.  Tuples, signed integers, rationals, arrays, finite sets, graphs, trees, matrices, formulas of a fixed grammar, and transcripts are represented by primitive-recursive codes.  There is no requirement that all applications share one global encoding.  A theorem-specific encoding is often smaller and easier to verify.

The checked repository already supplies a pairing function
\[
  \langle a,b\rangle=(a+b)^2+a
\]
with primitive-recursive projections and object proofs of the projection equations.  Repeated pairing is adequate for fixed arity.  For arrays whose length is a query parameter, a mixed-radix or bounded-sequence representation is more convenient.

\subsection{Mixed-radix arrays}

Let \(r_0,\ldots,r_{m-1}\geq 1\) be radices and let \(0\leq d_i<r_i\).  Define
\begin{equation}
  \code(d_0,\ldots,d_{m-1})
  =\sum_{i<m} d_i\prod_{j<i}r_j.
  \label{eq:mixed-code}
\end{equation}
The total range is
\[
  R=\prod_{i<m}r_i.
\]
The \(i\)-th digit is obtained by quotienting by \(\prod_{j<i}r_j\) and taking remainder modulo \(r_i\).  All operations are primitive recursive.  The checked division-with-remainder library supplies the arithmetic needed for the object proof.

\begin{lemma}[Mixed-radix correctness]
For every bounded digit vector, \(\decode_i(\code(\vec d))=d_i\).  Conversely, if \(c<R\), then re-encoding its decoded digits gives \(c\).
\end{lemma}

\begin{proof}
For the first direction, split the sum in \eqref{eq:mixed-code} below, at, and above \(i\).  Terms below \(i\) form a number smaller than the place value \(P_i=\prod_{j<i}r_j\).  The term at \(i\) is \(d_iP_i\).  Every term above \(i\) is divisible by \(P_ir_i\).  Division by \(P_i\), followed by remainder modulo \(r_i\), therefore leaves exactly \(d_i\).

For the converse, repeatedly apply the quotient-remainder equation
\[
  c=r_0(c\mathbin{\mathrm{quo}}r_0)+(c\mathbin{\mathrm{rem}}r_0).
\]
After \(m\) steps the remaining quotient is zero because \(c<R\).  Substitution yields the expansion \eqref{eq:mixed-code}.  In \iPRA, the variable-length statement is proved by quantifier-free induction on \(m\), with the current quotient and place value carried in the state.
\end{proof}

Mixed radix makes a master certificate literal: each node receives one bounded digit block, and all blocks form one number below a primitive-recursive product.  It also makes certificate minimality meaningful, because the search order is fixed.

\subsection{Finite predicates and characteristic bits}

A predicate used inside a loop is represented by a normalized bit \(b\) satisfying
\begin{equation}
  b=0\;\vee\;b=1.
  \label{eq:bit-normal}
\end{equation}
The repository's equality bit \texttt{eqb} and order development provide the atomic cases.  Boolean connectives are fixed primitive-recursive operations, for example
\[
\begin{aligned}
  \mathsf{not}(b)&=1\mathbin{\dot{-}} b,\\
  \mathsf{and}(b,c)&=bc,\\
  \mathsf{or}(b,c)&=\mathsf{sg}(b+c).
\end{aligned}
\]
Because inputs are normalized, their truth tables are derivable by finite case distinction in the intuitionistic object theory.

\subsection{Local proof contracts}

\begin{definition}[Finite contract]
A finite contract for queries \(q\) and outputs \(y\) consists of:
\begin{enumerate}[label=(\roman*),leftmargin=2.3em]
\item a certificate bound \(M(q)\);
\item a Boolean checker \(C(q,c)\);
\item an output decoder \(O(q,c)\);
\item residual formulas \(B_1(q),\ldots,B_r(q)\);
\item an object derivation
\[
  C(q,c)=1,\;B_1(q),\ldots,B_r(q),\Gamma(q)
  \;\vdash[\PRAAx]\;\Phi(q,O(q,c)).
\]
\end{enumerate}
All displayed programs are fixed primitive-recursive symbols after theorem specialization.
\end{definition}

A contract deliberately separates \emph{soundness} from \emph{inhabitation}.  The checker may be perfectly sound even when no accepted certificate exists.  Inhabitation is either derived by a producer theorem or retained as the residual
\begin{equation}
  \exists c<M(q)\;[C(q,c)=1].
  \label{eq:certificate-exists}
\end{equation}
This distinction is essential in Diophantine geometry, effective analysis, and compactness arguments: a finite verifier does not create a height bound, modulus, or cover.

\subsection{Composition of contracts}

Suppose contract \(A\) produces \(y\) from \(q\), and contract \(B\) produces \(z\) from \((q,y)\).  A composite master certificate is the pair \(\langle c_A,c_B\rangle\), or the corresponding mixed-radix block.  Its checker is
\[
 C_A(q,c_A)\wedge C_B(q,O_A(q,c_A),c_B),
\]
and its output is \(O_B(q,O_A(q,c_A),c_B)\).  The proof is a substitution of the first soundness theorem into the second, followed by simultaneous cut.  Residual lists concatenate, with duplicates optionally removed by a metalevel normalization.

No categorical machinery is needed for this basic operation, though finite contracts do form a useful category at the metalevel.  The crucial point is that composition preserves the object-theory endpoint and exposes, rather than hides, the union of residual producers.
